%   Modelo de trabalho completo ou resumo expandido para a XII Jornada Científica do IFRJ campus Volta Redonda
%   Em caso de dúvidas, enviar e-mail para copi.cvr@ifrj.edu.br
%   Codificação utilizada: UTF-8
%   Tamanho da tabulação: 4 (espaços)

\documentclass{jornada} % Classe do documento
%------------------------------------------------------
% ALGUNS PACOTES UTILIZADOS 
%------------------------------------------------------
\usepackage{enumitem,lipsum}

%--------------------------------------------------
% DADOS DO TRABALHO  
%--------------------------------------------------
\titulo{Título do resumo expandido submetido à XIV Jornada Científica}
\author[1*]{Nome do autor A} % Usar * para identificar o autor correspondente
\author[2]{Nome do autor B}
\author[1]{Nome do autor C}
\affil[1]{Instituição dos autores A e C (colocar o nome da instituição sem abreviações)}
\affil[2]{Instituição do autor B}
\email{email@.????} % colocar o e-amil do autor correspondente
\GrupoTematico{nome do grupo temático} % Inserir o grupo temático


% ARQUIVO DAS REFERÊNCIAS BIBLIOGRÁFICAS
\addbibresource{refbase.bib}

%------------------------------------------------------
% INÍCIO DO DOCUMENTO                                           %------------------------------------------------------        
\begin{document}
\pretext % não apague esta linha

% INÍCIO DA PARTE TEXTUAL                                           

\section{Introdução}

Lorem ipsum dolor sit amet, consectetuer adipiscing elit. Ut purus elit, vestibulum ut, placerat ac,adipiscing vitae, felis. Curabitur dictum gravida mauris. Nam arcu libero, nonummy eget, consectetuerid, vulputate a, magna. Donec vehicula augue eu neque. Pellentesque habitant morbi tristique senectuset netus et malesuada fames ac turpis egestas.  Mauris ut leo.  Cras viverra metus rhoncus sem.  Nullaet lectus vestibulum urna fringilla ultrices. Phasellus eu tellus sit amet tortor gravida placerat. Integersapien est, iaculis in, pretium quis, viverra ac, nunc. Praesent eget sem vel leo ultrices bibendum. Aeneanfaucibus. Morbi dolor nulla, malesuada eu, pulvinar at, mollis ac, nulla. Curabitur auctor semper nulla. Donec varius orci eget risus. Duis nibh mi, congue eu, accumsan eleifend, sagittis quis, diam. Duis egetorci sit amet orci dignissim rutrum.

\begin{subequations}
\label{eq:Maxwell}
Equações de Maxwell:
\begin{align}
        B'&=-\nabla \times E,         \label{eq:MaxB} \\
        E'&=\nabla \times B - 4\pi j, \label{eq:MaxE}
\end{align}
\end{subequations}

Nam dui ligula, fringilla a, euismod sodales, sollicitudin vel, wisi. Morbi auctor lorem non justo\footnote{Exemplo de nota de rodapé }. Nam lacus libero, pretium at, lobortis vitae, ultricies et, tellus.  Donec aliquet, tortor sed accumsanbibendum, erat ligula aliquet magna, vitae ornare odio metus a mi. Morbi ac orci et nisl hendrerit mollis.Suspendisse ut massa. Cras nec ante. Pellentesque a nulla\footnote{Exemplo de mais uma nota de rodapé com um texto bem mais longo de modo que ocupe mais que uma linha Cum sociis natoque penatibus et magnis disparturient montes, nascetur ridiculus mus. Aliquam tincidunt urna.Cum sociis natoque penatibus et magnis disparturient montes, nascetur ridiculus mus. Aliquam tincidunt urna.Cum sociis natoque penatibus et magnis disparturient montes, nascetur ridiculus mus. Aliquam tincidunt urna. Cum sociis natoque penatibus et magnis disparturient montes, nascetur ridiculus mus. Aliquam tincidunt urna.Cum sociis natoque penatibus et magnis disparturient montes, nascetur ridiculus mus. Aliquam tincidunt urna.Cum sociis natoque penatibus et magnis disparturient montes, nascetur ridiculus mus. Aliquam tincidunt urna.}. Cum sociis natoque penatibus et magnis disparturient montes, nascetur ridiculus mus. Aliquam tincidunt urna. Nulla ullamcorper vestibulum turpis. 
\begin{citacao}
 Quisque ullamcorper placerat ipsum. Cras nibh. Morbi vel justo vitae lacus tinciduntultrices.  Lorem ipsum dolor sit amet, consectetuer adipiscing elit.  In hac habitasseplatea dictumst.  Integer tempus convallis augue.  Etiam facilisis.  Nunc elementumfermentum wisi.  Aenean placerat.  Ut imperdiet, enim sed gravida sollicitudin, felisodio placerat quam, ac pulvinar elit purus eget enim. Nunc vitae tortor. Proin tempusnibh sit amet nisl. Vivamus quis tortor vitae risus porta vehicula.
\end{citacao}

\lipsum[5]

\section{Elementos flutuantes}
Tabelas, quadros e figuras são elementos utilizados em trabalhos acadêmicos, e possuem diferenças e especificações definidas pela ABNT.

\subsection{Exemplo se subseção}
As tabelas são formadas por linhas verticais, devem manter suas bordas laterais abertas e geralmente são utilizadas para dados quantitativos. Os quadros, por outro lado, são formados por linhas verticais e horizontais, devem ter todas suas extremidades fechadas e são mais utilizados para dados qualitativos.

Por último, as figuras são elementos ilustrativos, que podem ser em forma de fotos, mapas, gráficos, gravuras, etc.

\begin{table}[H]
\IBGEtab{
  \caption{Pessoas Residentes em domicílios particulares, por sexo e situação do domicílio - Brasil - 1980}%
  \label{tab: residentes}
  }
{
\begin{tabular}{m{2.0cm}ccc}
 \toprule
 Situação do Total & Total & Mulheres & Homens \\
\midrule
Total & 117.960.301 & 59.595.332 & 58.364.969 \\
Urbana & 79.972.931 & 41.115.439 & 38.857.492 \\
Rural & 37.987.370 & 18.479.893 & 19.507.477 \\
\bottomrule
\end{tabular}
}
{
\fonte{Elaborada pelo autor}
\nota{Dados fictícios}
}
\end{table}
		
		
\begin{quadro}[H]
	\centering
	\caption{Alguns detalhes sobre elementos flutuantes}
	\label{quadroex1}
	\begin{tabular}{|l|m{4cm}|m{4cm}|m{4cm}|}
		\hline
		& \textbf{Tabela}                                                                             & \textbf{Quadro}                                                                             & \textbf{Figuras}                                                                    \\ \hline
		\textbf{Formato}    & Bordas laterais não podem ser fechadas.                                                     & As extremidades devem ser fechadas.                                                         & Podem ser em forma de fotos, mapas, gráficos, gravuras, etc.                        \\ \hline
		\textbf{Uso}        & Geralmente para dados quantitativos.                                                        & Geralmente para dados qualitativos.                                                         & Ilustrar informações e dados.                                                       \\ \hline
		\textbf{Elementos}  & Título, cabeçalho, conteúdo, fonte e, se necessário, notas explicativas.                    & Título, fonte, legenda e notas.                                                             & Título, numeração e fonte.                                                          \\ \hline
		\textbf{Divisão}    & Formada por linhas horizontais.                                                               & Formado por linhas horizontais e verticais.                                                 & -                                                                                   \\ \hline
		\textbf{Formatação} & O número e o título da tabela devem vir acima dela, enquanto a fonte deve aparecer embaixo. & O número e o título do quadro devem vir acima dele, enquanto a fonte deve aparecer embaixo. & O número e o título devem aparecer no topo, enquanto a fonte deve aparecer embaixo. \\ \hline
	\end{tabular}
	\fonte{Elaborado pelo autor}
\end{quadro}
				
Abaixo, a figura \ref{fig:ex} é apresentada como exemplo de disposição dos elementos legenda, numeração e fonte.
Para uso avançado de imagens no \LaTeX, recomenda-se a consulta de literatura especializada \cite{Goossens2007}.

	\begin{figure}[H]
		\centering
		\caption{Exemplo de inclusão de imagem}
		\includegraphics[height=6cm]{./figuras/wave.png}
		\fonte{\textcite{opensource}}
		\label{fig:ex}
	\end{figure}

\subsection{Outros exemplos}
Nesta subseção apresentamos alguns elementos deste modelo em \LaTeX.

\subsubsection{Citações diretas com mais de três linhas}
Neste modelo, para fazermos uma citação direta com mais de três linhas, podemos usar o ambiente \texttt{citacao}. Por exemplo, se digitarmos
\begin{verbatim}
\begin{citacao}
    Lorem ipsum dolor sit amet, consectetur adipiscing elit.
    Nulla in pulvinar nulla. Sed tempor condimentum felis quis blandit.
    Maecenas ultricies sapien eu nisi hendrerit maximus.
    Nullam sagittis consectetur leo, at feugiat lorem.
    Nulla sed maximus metus, sed tincidunt leo.
    Fusce fermentum consectetur libero, eget rhoncus nisl laoreet eget.
    Nam eget justo metus.
\end{citacao}
\end{verbatim}
obteremos o seguinte resultado
\begin{citacao}
    Lorem ipsum dolor sit amet, consectetur adipiscing elit. Nulla in pulvinar nulla. Sed tempor condimentum felis quis blandit. Maecenas ultricies sapien eu nisi hendrerit maximus. Nullam sagittis consectetur leo, at feugiat lorem. Nulla sed maximus metus, sed tincidunt leo. Fusce fermentum consectetur libero, eget rhoncus nisl laoreet eget. Nam eget justo metus.
\end{citacao}

Lorem ipsum dolor sit amet, consectetur adipiscing elit. Nulla in pulvinar nulla. Sed tempor condimentum felis quis blandit.
%----------------------------
% Parte pós-textual

\section*{Agradecimentos}
Esta é uma seção reservada aos agradecimentos dos autores, caso pertinente.	

% REFERÊNCIAS BIBLIOGRÁFICAS

\singlespacing
\printbibliography
\end{document}